% GAP 1 + MOONSHOT: High-Dimensional Fourier Stability + Gardner Problem 1.5
% Replaces the planar-only Theorem 4 with a general-D result via optimal recovery theory.
% Also states partial progress on Problem 1.5.

% ===================================================================
% KEY INSIGHT (supersedes the planar proof)
% ===================================================================
%
% The planar proof (theorem4_proof.tex) used Jackson's theorem + trigonometric
% interpolation on S^1. This works beautifully for D=2 but the naive extension
% to S^{D-1} via spherical harmonics is killed by the Lebesgue constant
% (grows polynomially in the degree, consuming the approximation rate).
%
% The FIX: use optimal recovery theory, which gives the minimax rate directly
% without going through explicit interpolation. The rate depends only on the
% smoothness class and the dimension of the sphere.
%
% For C^2 functions on S^{D-1}: optimal recovery rate from m function values
% is m^{-2/(D-1)}. This is the EXACT rate of Theorem 4.

% ===================================================================
% THEOREM 4: GENERAL DIMENSION
% ===================================================================

\begin{theorem}[Stability of Convex Body Determination from Width Measurements]
\label{thm:stability-general}
Let $K, L \subset B_R^D$ be origin-symmetric convex bodies with Gauss curvature bounded below by $\kappa > 0$ (equivalently, the principal radii of curvature satisfy $r_i \geq \kappa$ for all $i$). Let their support functions $h_K, h_L \in C^2(S^{D-1})$, and suppose that $m$ width measurements agree within $\varepsilon$:
\[
  |w_K(u_i) - w_L(u_i)| \leq \varepsilon \quad \text{for } i = 1, \ldots, m
\]
where $\{u_i\}_{i=1}^m$ are $m$ measurement directions on $S^{D-1}$.

Then the Hausdorff distance satisfies:
\[
  \boxed{\delta_H(K, L) \leq C_D \!\left(\varepsilon + \frac{R}{\kappa \, m^{2/(D-1)}}\right)}
\]
where $C_D$ depends only on $D$. Moreover, this rate is minimax optimal: there exist pairs of convex bodies achieving $\delta_H \geq c_D \cdot R / (\kappa \, m^{2/(D-1)})$ for any choice of $m$ measurement directions.
\end{theorem}

\begin{proof}
The proof has three components: the upper bound, the lower bound (tightness), and the noise term.

% --- UPPER BOUND ---
\textbf{Upper Bound.}
Let $g = h_K - h_L : S^{D-1} \to \mathbb{R}$. Since $K$ and $L$ have curvature $\geq \kappa$, the support functions satisfy $h_K, h_L \in C^2(S^{D-1})$ with:
\[
  \|h_K\|_{C^2(S^{D-1})}, \|h_L\|_{C^2(S^{D-1})} \leq C \cdot R / \kappa.
\]
This follows from the relationship between support function regularity and curvature: the eigenvalues of $\nabla^2 h_K + h_K \cdot g_{S^{D-1}}$ (where $g_{S^{D-1}}$ is the round metric) equal the principal radii of curvature of $\partial K$.

The key tool is the theory of \emph{optimal recovery} \cite{micchelli1977survey, traub1988information}. For a function class $\mathcal{F} \subset C(S^{D-1})$ and $m$ point evaluations at locations $\{u_i\}$, the \emph{optimal recovery error} is:
\[
  E_m^*(\mathcal{F}) = \inf_{\{u_i\}, \phi} \sup_{f \in \mathcal{F}} \|f - \phi(f(u_1), \ldots, f(u_m))\|_\infty
\]
where the infimum is over all choices of sample points and all recovery maps $\phi : \mathbb{R}^m \to C(S^{D-1})$.

For the Sobolev ball $\mathcal{F} = W^{2,\infty}(S^{D-1}) \cap B_M$ (functions with bounded $C^2$ norm, with $M = CR/\kappa$), the optimal recovery rate is:
\[
  E_m^*(\mathcal{F}) \asymp M \cdot m^{-2/(D-1)}.
\]

This follows from two classical results:

\emph{(i) Kolmogorov $n$-widths.} The $n$-th Kolmogorov width of $W^{2,\infty}(S^{D-1}) \cap B_M$ in $C(S^{D-1})$ satisfies:
\[
  d_n(W^{2,\infty}(S^{D-1}) \cap B_M, C(S^{D-1})) \asymp M \cdot n^{-2/(D-1)}.
\]
The upper bound comes from spherical polynomial approximation of degree $L$, which uses $n \sim L^{D-1}$ basis functions, and Jackson's inequality on $S^{D-1}$ \cite{ragozin1970polynomial, dai2013approximation}:
\[
  E_L(f)_\infty \leq C_D \frac{\omega_2(f, 1/L)_\infty}{1} \leq \frac{C_D}{L^2} \|f\|_{C^2}
\]
where $\omega_2$ is the second modulus of smoothness on $S^{D-1}$.

The lower bound is the Bernstein inequality: there exist trigonometric/spherical polynomial approximation barriers matching the Jackson rate \cite{ditzian1990moduli}.

\emph{(ii) Equivalence of widths and optimal recovery.} For linear problems on convex symmetric function classes, the optimal recovery error from $m$ function values is equivalent (up to constants) to the $m$-th Kolmogorov width \cite{magaril1979diameters, magaril2006optimal}. Formally:
\[
  E_m^*(\mathcal{F}) \asymp d_m(\mathcal{F}, C(S^{D-1})) \asymp M \cdot m^{-2/(D-1)}.
\]

Since $g = h_K - h_L \in W^{2,\infty}(S^{D-1})$ with $\|g\|_{C^2} \leq 2CR/\kappa$, the optimal recovery from $m$ noiseless samples gives:
\[
  \|g\|_\infty \leq C_D \cdot \frac{R}{\kappa} \cdot m^{-2/(D-1)}.
\]
And $\delta_H(K, L) = \|h_K - h_L\|_\infty = \|g\|_\infty$.

% --- NOISE ---
\textbf{Noise term.} When samples are corrupted by noise $\varepsilon$: $|g(u_i)| \leq \varepsilon$ for all $i$, the recovery error includes both the approximation error (from finite sampling) and the measurement error. The measurement error propagates through the optimal recovery map at rate $O(\varepsilon)$ for well-conditioned recovery (which is guaranteed when sample points are quasi-uniform). Thus:
\[
  \delta_H(K, L) \leq C_D\!\left(\varepsilon + \frac{R}{\kappa \, m^{2/(D-1)}}\right).
\]

% --- LOWER BOUND ---
\textbf{Lower Bound (Tightness).}
We must show that for any $m$ directions $\{u_i\}$, there exist convex bodies $K, L$ with curvature $\geq \kappa$ and $w_K(u_i) = w_L(u_i)$ for all $i$, yet $\delta_H(K, L) \geq c_D \cdot R/(\kappa \, m^{2/(D-1)})$.

Construct $g$ as a spherical polynomial vanishing at all sample points. The space of spherical polynomials of degree $\leq 2L$ (where $L = \lceil m^{1/(D-1)} \rceil$) has dimension $\dim \mathcal{P}_{2L}(S^{D-1}) \asymp (2L)^{D-1} \asymp 2^{D-1} m$. Imposing the $m$ linear conditions $g(u_i) = 0$ leaves a subspace of dimension $\geq (2^{D-1} - 1)m > 0$. Choose $g$ in this null space with maximal $L^\infty$ norm subject to $\|g\|_{C^2} \leq CR/\kappa$.

By the Bernstein inequality on $S^{D-1}$ \cite{dai2013approximation}: for a spherical polynomial of degree $\leq N$, $\|\nabla^2 p\|_\infty \leq C N^2 \|p\|_\infty$. Applied to $g$ with $N = 2L$:
\[
  \|g\|_{C^2} \asymp (2L)^2 \|g\|_\infty \leq CR/\kappa \implies \|g\|_\infty \leq \frac{CR}{4\kappa L^2}.
\]
Since $g$ is in the null space, $g(u_i) = 0$ for all $i$, so $K' = \{x : h_{K'}(u) = h_K(u) + g(u)\}$ is $(\varepsilon, \Pi)$-indistinguishable from $K$ (with $\varepsilon = 0$). The convexity of $K'$ is guaranteed when $\|g\|_{C^2}$ is small relative to $\kappa$ (the perturbation $g$ changes the curvature by at most $\|g\|_{C^2}$, so curvature $\geq \kappa - CR/\kappa > 0$ for $C$ small). Then:
\[
  \delta_H(K, K') = \|g\|_\infty \geq c \cdot \frac{R}{\kappa L^2} \asymp \frac{R}{\kappa \, m^{2/(D-1)}}.
\]
This matches the upper bound rate. \qed
\end{proof}

\begin{corollary}[Rate Comparison: Smooth vs.\ Lipschitz]
\label{cor:rates}
The smooth stability bound (Theorem~\ref{thm:stability-general}) and the Lipschitz bound (Theorem~\ref{thm:indist-corrected}) compare as:

\begin{center}
\renewcommand{\arraystretch}{1.3}
\begin{tabular}{ccccc}
\toprule
$d_{\mathrm{eff}}$ & \textbf{Lipschitz rate} & \textbf{Smooth rate} & \textbf{Improvement} & \textbf{Benchmarks to halve gap} \\
\midrule
2 & $m^{-1}$ & $m^{-2}$ & $m^{-1}$ & $\sqrt{2}\times$ vs $2\times$ \\
3 & $m^{-1/2}$ & $m^{-1}$ & $m^{-1/2}$ & $\sqrt{2}\times$ vs $4\times$ \\
4 & $m^{-1/3}$ & $m^{-2/3}$ & $m^{-1/3}$ & $2^{3/2}\times$ vs $8\times$ \\
5 & $m^{-1/4}$ & $m^{-1/2}$ & $m^{-1/4}$ & $4\times$ vs $16\times$ \\
$D$ & $m^{-1/(D-1)}$ & $m^{-2/(D-1)}$ & $m^{-1/(D-1)}$ & $2^{(D-1)/2}\times$ vs $2^{D-1}\times$ \\
\bottomrule
\end{tabular}
\end{center}

\noindent In all dimensions, the smooth rate \emph{squares} the Lipschitz rate. Equivalently: for a target accuracy $\delta$, the smooth bound requires $\sqrt{m_{\mathrm{Lip}}}$ benchmarks where $m_{\mathrm{Lip}}$ is the Lipschitz requirement.
\end{corollary}

\begin{proof}
Immediate from $m^{-2/(D-1)} = (m^{-1/(D-1)})^2$ and inverting: $m_{\mathrm{smooth}} = (CR/(\kappa\delta))^{(D-1)/2}$ vs $m_{\mathrm{Lip}} = (CR/\delta)^{D-1}$, so $m_{\mathrm{smooth}} = m_{\mathrm{Lip}}^{1/2} \cdot (\kappa \text{ correction})$. \qed
\end{proof}


% ===================================================================
% MOONSHOT: TOWARD GARDNER'S PROBLEM 1.5
% ===================================================================

\subsection{Connection to Gardner's Problem 1.5}

Gardner \cite{gardner1995geometric} posed the following:

\begin{problem}[Gardner 1.5]
How stably does a finite number of X-rays determine a convex body?
\end{problem}

An \emph{X-ray} of $K$ in direction $u$ is the function:
\[
  X_u K(x) = \lambda_1(K \cap (x + \mathbb{R}u)), \quad x \in u^\perp,
\]
giving the chord length of $K$ at each point $x$ in the hyperplane perpendicular to $u$.

\textbf{Relationship to our setting.} Width measurements are \emph{one-dimensional summaries} of X-rays:
\[
  w_K(u) = \int_{u^\perp} X_u K(x)\, d\lambda_{D-1}(x) \cdot \frac{1}{\text{vol}_{D-1}(\Pi_{u^\perp}(K))}... 
\]
More precisely, the width is related to the support function, while the X-ray contains strictly more information (the full chord-length profile, not just the total).

\begin{proposition}[Width stability as a lower bound for X-ray stability]
\label{prop:xray-lb}
For any convex bodies $K, L$ and any set of $m$ directions, if the X-rays agree within $\varepsilon$ (in $L^1(u^\perp)$ norm for each direction):
\[
  \|X_{u_i} K - X_{u_i} L\|_{L^1(u_i^\perp)} \leq \varepsilon' \quad \text{for all } i,
\]
then the widths agree within $\varepsilon = C \varepsilon'$. Consequently, any lower bound on recovery from widths is also a lower bound on recovery from X-rays:
\[
  E_m^*(\text{widths}) \leq E_m^*(\text{X-rays}).
\]
\end{proposition}

\begin{proof}
The width $w_K(u) = h_K(u) + h_K(-u)$ satisfies:
\[
  |w_K(u) - w_L(u)| \leq \|X_u K - X_u L\|_{L^1(u^\perp)} \leq \varepsilon'
\]
since the width is the integral of the chord-length function. The bound on $E_m^*$ follows from the fact that X-rays provide strictly more information than widths. \qed
\end{proof}

\begin{conjecture}[Partial Answer to Problem 1.5]
\label{conj:problem15}
For origin-symmetric convex bodies in $\mathbb{R}^D$ with curvature $\geq \kappa$:

\begin{enumerate}[label=(\alph*)]
  \item \textbf{Width measurements:} $\delta_H(K, L) \leq C_{D,\kappa,R} \cdot m^{-2/(D-1)}$ (Theorem~\ref{thm:stability-general}, proven).
  
  \item \textbf{X-ray measurements:} We conjecture $\delta_H(K, L) \leq C_{D,\kappa,R} \cdot m^{-\alpha(D)}$ where $\alpha(D) > 2/(D-1)$, reflecting the additional information in X-rays beyond widths.
  
  \item \textbf{For $D = 2$:} $\alpha(2) = 2$ for widths (proven). For X-rays, the rate should be at least $\alpha = 2$, and may be faster due to the Fourier slice theorem: $m$ X-ray directions determine $m$ lines in Fourier space, giving $O(m^2)$ effective constraints.
\end{enumerate}
\end{conjecture}

\begin{remark}[What this means for LLM evaluation]
Benchmarks are scalar measurements --- closer to widths than X-rays. Thus Theorem~\ref{thm:stability-general} is the directly applicable result. However, richer evaluation methods (e.g., full response distributions rather than accuracy scores, or item-level results rather than aggregate scores) correspond to X-ray-like measurements and should enable substantially better discrimination between models. This connects to the IRT (Item Response Theory) approach of \citet{polo2024tinybenchmarks}: item-level analysis gives ``X-ray'' information, while aggregate scores give only ``width'' information.
\end{remark}


% ===================================================================
% THE PRACTICAL MESSAGE
% ===================================================================

\subsection{Practical Implications}

\begin{enumerate}
  \item \textbf{The Lipschitz bound is not tight.} Convex bodies with bounded curvature have C$^2$ support functions, and the optimal recovery rate is $m^{-2/(D-1)}$, squaring the Lipschitz rate $m^{-1/(D-1)}$. This means the blind spot shrinks faster with well-designed benchmarks than the Lipschitz bound suggests.
  
  \item \textbf{But the improvement requires curvature.} The $m^{-2/(D-1)}$ rate holds for convex bodies with curvature bounded away from zero. For bodies with flat faces (curvature $= 0$), the Lipschitz rate $m^{-1/(D-1)}$ is tight. In LLM terms: if capabilities are ``smooth'' (no sharp thresholds), the blind spot shrinks faster.
  
  \item \textbf{The square-root law.} Theorem 4 says: whatever the Lipschitz bound says you need, the smooth bound says you need only the \emph{square root} as many benchmarks. For $d_{\mathrm{eff}} = 5$ and a target accuracy of $\delta = 0.1R$: Lipschitz needs $m \geq (10C)^4 \approx 10{,}000$ benchmarks; smooth needs $m \geq (10C)^2 \approx 100$ benchmarks. The gap is dramatic.
  
  \item \textbf{Connection to Theorem 3.} The greedy algorithm (Theorem 3) should prioritize benchmarks that improve the \emph{covering quality} of directions on $S^{d_{\mathrm{eff}}-1}$, not just maximize marginal variance. Combining the greedy algorithm with the stability bound gives a principled answer to the ``next benchmark'' question.
\end{enumerate}
